\section*{Chapter \ref{ch:b3:lebesgue} --- The Lebesgue
Integral}

\begin{solution}{exo:b3:lebesgue:1}
(a) If $f$ is nondecreasing, $\{f > t\}$ is $\varnothing$, $\R$,
or a ray $\intoo a{+\infty}$ / $\intco a{+\infty}$: Borel in
every case; nonincreasing likewise. A derivative: $f'(x) =
\lim_n n\bigl(f(x + \frac1n) - f(x)\bigr)$ is a pointwise limit
of continuous (hence measurable) functions:
\cref{prop:b3:lebesgue:stability}(c).

(b) $\{f > t\} = \bigcup_{q \in \Q,\, q > t}\{f > q\}$: if the
rational levels are measurable, all levels are, and the rays
generate $\mathcal B(\R)$.
\end{solution}

\begin{solution}{exo:b3:lebesgue:2}
First: $(1 + x/n)^n \nearrow \eu^x$ for $x \geq 0$, so the
integrand tends pointwise to $\eu^{-x}\cos x$; for $n \geq 2$,
$(1 + x/n)^n \geq (1 + x/2)^2$, giving the integrable dominator
$(1 + x/2)^{-2}$. DCT:
\[
\lim_n\int_0^\infty\frac{\cos x}{(1 + x/n)^n}\dd x
= \int_0^\infty \eu^{-x}\cos x\,\dd x
= \operatorname{Re}\int_0^\infty\eu^{-(1 - \iu)x}\dd x
= \operatorname{Re}\frac{1}{1 - \iu} = \frac12 .
\]
Second: substitute $u = x^n$ (a $\mathcal C^1$ bijection of
$\intoo01$):
\[
\int_0^1\frac{nx^{n-1}}{1 + x}\dd x = \int_0^1\frac{\dd u}{1 +
u^{1/n}} \longrightarrow \int_0^1\frac{\dd u}{2} = \frac12,
\]
by DCT: for $u \in \intoo01$, $u^{1/n} \to 1$, and the
integrand is bounded by $1$ on a finite measure space.
\end{solution}

\begin{solution}{exo:b3:lebesgue:3}
(a) $f_n = n\,\mathbf 1_{\intoo0{1/n}}$: $\liminf f_n = 0$
pointwise, $\int f_n = 1$: $0 < 1$.

(b) Height: $n\mathbf 1_{\intoo0{1/n}}$; width: $\frac1n\mathbf
1_{\intoo0n}$; translation: $\mathbf 1_{\intoo n{n+1}}$. All
tend to $0$ pointwise with $\int = 1$. In each case the
\emph{domination} hypothesis fails: $\sup_nf_n$ is
$\approx 1/x$ near $0$, $\approx$ a nonintegrable constant
profile, $\mathbf 1_{\intoo1\infty}$-like --- never integrable.

(c) ``$\int\limsup f_n \geq \limsup\int f_n$'' fails for the
translating bump: left side $0$, right side $1$.
``$\limsup\int \leq \int\limsup$'' is the same statement. And
Fatou for $\limsup$ with $\leq$ reversed
(``reverse Fatou'') requires a dominator --- the same bump is
the counterexample.
\end{solution}

\begin{solution}{exo:b3:lebesgue:4}
(a) For $x > 0$: $\frac1{\eu^x - 1} = \frac{\eu^{-x}}{1 -
\eu^{-x}} = \sum_{n\geq1}\eu^{-nx}$, so $\frac{x}{\eu^x - 1} =
\sum_{n\geq1}x\eu^{-nx}$, a series of nonnegative measurable
functions: \cref{cor:b3:lebesgue:additivity} allows term-by-term
integration:
\[
\int_0^\infty\frac{x\,\dd x}{\eu^x - 1}
= \sum_{n\geq1}\int_0^\infty x\eu^{-nx}\dd x
= \sum_{n\geq1}\frac1{n^2} = \frac{\pi^2}6
\]
($\int_0^\infty x\eu^{-nx}\dd x = n^{-2}$ by parts; Basel from
the Year 2 volume, or \cref{exo:b3:hilbert:5} to come).

(b) On $\intoo01$, $-x\ln x \geq 0$, so $x^{-x} =
\eu^{-x\ln x} = \sum_k\frac{(-x\ln x)^k}{k!}$ is a series of
nonnegative terms: interchange again. Substituting $x =
\eu^{-u/(k+1)}$:
\[
\int_0^1(-x\ln x)^k\dd x
= \int_0^\infty\Bigl(\frac{u}{k+1}\Bigr)^{k}
\eu^{-\frac{ku}{k+1}}\;\frac{\eu^{-\frac u{k+1}}}{k+1}\,\dd u
= \frac{1}{(k+1)^{k+1}}\int_0^\infty u^k\eu^{-u}\dd u
= \frac{k!}{(k+1)^{k+1}} .
\]
Hence $\int_0^1x^{-x}\dd x = \sum_{k\geq0}\frac{1}{(k+1)^{k+1}}
= \sum_{n\geq1}n^{-n}$: the sophomore's dream, rigorously.
\end{solution}

\begin{solution}{exo:b3:lebesgue:5}
(a) Markov: $a\,\mathbf 1_{\{f \geq a\}} \leq f$, integrate:
$\mu(\{f \geq a\}) \leq \frac1a\int f$. If $\int f = 0$:
$\mu(\{f \geq \frac1n\}) = 0$ for every $n$, and $\{f > 0\} =
\bigcup_n\{f \geq \frac1n\}$ is null.

(b) $\mu(\{\abs f = \infty\}) \leq \mu(\{\abs f \geq n\}) \leq
\frac1n\int\abs f \to 0$.

(c) Take $A = \{f > 0\}$: $\int f^+\dd\mu = \int_Af\,\dd\mu =
0$, so $f^+ = 0$ a.e.\ by (a); likewise $f^- = 0$ a.e.
\end{solution}

\begin{solution}{exo:b3:lebesgue:6}
(a) $\mathbf 1_\Q$: discontinuous everywhere, not
Riemann-integrable (\cref{thm:b3:lebesgue:riemann}); its
Lebesgue integral is $\lambda(\Q) = 0$. Thomae: continuous at
every irrational (given $\varepsilon$, only finitely many
rationals in $\intcc01$ have denominator $\leq
1/\varepsilon$; avoid them by a small neighborhood),
discontinuous at rationals (density of irrationals): continuous
a.e., Riemann-integrable, integral $0$ (it vanishes a.e.).
$\mathbf 1_K$, $K$ a fat Cantor set: the discontinuity set is
$\partial K = K$ (closed with empty interior), of measure
$\frac12 > 0$: \emph{not} Riemann-integrable --- yet
Lebesgue-integrable with integral $\lambda(K) = \frac12$.

(b) $\int_{k\pi}^{(k+1)\pi}\frac{\abs{\sin x}}x\dd x \geq
\frac1{(k+1)\pi}\int_{k\pi}^{(k+1)\pi}\abs{\sin x}\dd x =
\frac{2}{(k+1)\pi}$: the series diverges. Convergence of the
improper integral: for $A > \pi$,
\[
\int_\pi^A\frac{\sin x}x\dd x = \Bigl[\frac{-\cos
x}x\Bigr]_\pi^A - \int_\pi^A\frac{\cos x}{x^2}\dd x,
\]
and both terms converge as $A \to \infty$ ($\frac1{x^2}$ is
integrable): conditional convergence without absolute
integrability.
\end{solution}

\begin{solution}{exo:b3:lebesgue:7}
Domination: $\abs{\partial_t(\eu^{-x^2}\cos(tx))} =
\abs{x\eu^{-x^2}\sin(tx)} \leq x\eu^{-x^2}$, integrable and
independent of $t$: \cref{thm:b3:lebesgue:paramdiff} applies
globally,
\[
F'(t) = -\int_0^\infty x\eu^{-x^2}\sin(tx)\,\dd x
= \Bigl[\tfrac12\eu^{-x^2}\sin(tx)\Bigr]_0^\infty
- \frac t2\int_0^\infty\eu^{-x^2}\cos(tx)\dd x
= -\frac t2F(t)
\]
(integration by parts with $\dd v = x\eu^{-x^2}\dd x$). The
linear ODE gives $F(t) = F(0)\eu^{-t^2/4}$; with $F(0) =
\frac{\sqrt\pi}2$ (\cref{pb:b3:lebesgue:1}), the Gaussian
reproduces itself under this cosine transform.
\end{solution}

\begin{solution}{exo:b3:lebesgue:8}
The integral converges: near $0$ the integrand tends to $b - a$
(bounded), and it decays like $\eu^{-ax}$ at infinity. Fix $a$
and view $I(b) = \int_0^\infty\frac{\eu^{-ax} -
\eu^{-bx}}x\dd x$ as a function of $b \in \intco a{+\infty}$.
For all $b \geq a$: $\abs{\partial_b(\text{integrand})} =
\eu^{-bx} \leq \eu^{-ax}$, and $\int_0^\infty\eu^{-ax}\dd x =
\frac1a < \infty$: an integrable dominator. So
\cref{thm:b3:lebesgue:paramdiff} gives $I'(b) =
\int_0^\infty\eu^{-bx}\dd x = \frac1b$, and $I(a) = 0$:
\[
I(b) = \int_a^b\frac{\dd t}t = \ln\frac ba .
\]
(Equivalently, the hint's route: the integrand is
$\int_a^b\eu^{-xt}\dd t \geq 0$ and the interchange is the
continuous analogue of \cref{cor:b3:lebesgue:additivity},
i.e.\ Tonelli --- proved in \cref{ch:b3:product}; the
parameter route stays within this chapter.)
\end{solution}

\begin{solution}{exo:b3:lebesgue:9}
$\nu(\varnothing) = 0$; for disjoint $(A_n)$, $f\mathbf
1_{\bigsqcup A_n} = \sum_nf\mathbf 1_{A_n}$ (pointwise, all
terms $\geq 0$), and \cref{cor:b3:lebesgue:additivity} gives
$\sigma$-additivity. The formula $\int g\,\dd\nu = \int
gf\,\dd\mu$: for $g = \mathbf 1_A$ it is the definition of
$\nu$; for simple $g$, linearity; for $g \geq 0$ measurable,
take simple $s_n \nearrow g$
(\cref{thm:b3:lebesgue:approximation}): $s_nf \nearrow gf$, and
MCT on both sides passes to the limit. (This ``indicator
$\to$ simple $\to$ MCT'' escalator is the standard machine of
the theory.)
\end{solution}

\begin{solution}{exo:b3:lebesgue:10}
(a) $\abs{\sin(tx)} \leq \abs tx$ gives
$\abs{\frac{\sin(tx)}{x(1+x^2)}} \leq \frac{\abs
t}{1+x^2}$: the integral converges, and on $\abs t \leq T$ the
dominator $\frac{T}{1+x^2}$ yields continuity
(\cref{thm:b3:lebesgue:paramcont}). Differentiation:
$\abs{\partial_t} = \abs{\frac{\cos(tx)}{1+x^2}} \leq
\frac1{1+x^2}$, integrable: $f'(t) =
\int_0^\infty\frac{\cos(tx)}{1+x^2}\dd x$ for all $t$. A second
differentiation would require integrating $\frac{x\sin(tx)}{1 +
x^2}$, whose absolute value behaves like $\frac{\abs{\sin(tx)}}
x$ at infinity: \emph{not} integrable --- no dominator exists
and \cref{thm:b3:lebesgue:paramdiff} cannot be applied again.

(b) Admitting $f'(t) = \frac\pi2\eu^{-t}$ for $t > 0$: by
oddness of $f$, $f'$ is even, so $f'(t) =
\frac\pi2\eu^{-\abs t}$ --- which is \emph{not differentiable
at $0$}: $f$ is $\mathcal C^1$ but not $\mathcal C^2$,
confirming that the blocked second differentiation was not a
technical accident. For $t > 0$ the improper integral
$\int_0^\infty\frac{x\sin(tx)}{1+x^2}\dd x$ equals $-f''(t) =
\frac\pi2\eu^{-t}$ (differentiating the admitted formula where
it is legitimate, i.e.\ on $\intoo0\infty$) --- an improper,
non-Lebesgue value.
\end{solution}

\begin{solution}{pb:b3:lebesgue:1}
\textbf{1.} $A$ is $\mathcal C^1$ by the fundamental theorem
of calculus and the chain rule: $A'(t) =
2\eu^{-t^2}\int_0^t\eu^{-x^2}\dd x$. For $B$:
$\partial_t\bigl(\frac{\eu^{-t^2(1+x^2)}}{1+x^2}\bigr) =
-2t\,\eu^{-t^2(1+x^2)}$, continuous and bounded on $[t_0, T]
\times \intcc01$ for any $0 < t_0 < T$ (bounded domination on a
finite measure space suffices): $B$ is $\mathcal C^1$ on
$\intoo0{+\infty}$ with
\[
B'(t) = -2t\int_0^1 \eu^{-t^2(1+x^2)}\dd x
= -2\eu^{-t^2}\int_0^1 t\,\eu^{-(tx)^2}\dd x
= -2\eu^{-t^2}\int_0^t\eu^{-u^2}\dd u = -A'(t)
\]
(substitution $u = tx$).

\textbf{2.} $A(0) = 0$ and $B(0) = \int_0^1\frac{\dd x}{1+x^2}
= \frac\pi4$. $A + B$ has zero derivative on
$\intoo0{+\infty}$ and is continuous at $0$ ($B$ by
domination $\frac1{1+x^2}$ and
\cref{thm:b3:lebesgue:paramcont}): $A + B \equiv \frac\pi4$.
As $t\to+\infty$: $0 \leq B(t) \leq \eu^{-t^2} \to 0$, and
$A(t) \to \bigl(\int_0^\infty\eu^{-x^2}\dd x\bigr)^2$ (MCT or
just monotone convergence of the inner integral).

\textbf{3.} Hence $\bigl(\int_0^\infty\eu^{-x^2}\dd x\bigr)^2
= \frac\pi4$: $\int_0^\infty\eu^{-x^2}\dd x =
\frac{\sqrt\pi}2$, and by evenness $G = \sqrt\pi$. Also
$\Gamma(\frac12) = \int_0^\infty x^{-1/2}\eu^{-x}\dd x
\overset{x = u^2}{=} 2\int_0^\infty\eu^{-u^2}\dd u =
\sqrt\pi$.

\textbf{4.} For $t > 0$: $\abs{\eu^{-tx}\frac{\sin x}x} \leq
\eu^{-tx}$, integrable. For $t = 0$, improper convergence: on
$[\pi, A]$,
\[
\int_\pi^A\frac{\sin x}x\dd x
= \Bigl[-\frac{\cos x}x\Bigr]_\pi^A -
\int_\pi^A\frac{\cos x}{x^2}\dd x,
\]
both terms convergent as $A \to \infty$; near $0$ the integrand
extends continuously by $1$.

\textbf{5.} On $[t_0, +\infty)$ ($t_0 > 0$):
$\abs{\partial_t\bigl(\eu^{-tx}\tfrac{\sin x}x\bigr)} =
\eu^{-tx}\abs{\sin x} \leq \eu^{-t_0x}$, integrable:
\cref{thm:b3:lebesgue:paramdiff} applies on every such
interval, so on all of $\intoo0{+\infty}$:
\[
F'(t) = -\int_0^\infty\eu^{-tx}\sin x\,\dd x
= -\operatorname{Im}\int_0^\infty\eu^{-(t - \iu)x}\dd x
= -\operatorname{Im}\frac{1}{t - \iu} = -\frac{1}{1 + t^2}.
\]

\textbf{6.} $\abs{F(t)} \leq \int_0^\infty\eu^{-tx}\dd x =
\frac1t \to 0$. Integrating $F' = -\frac1{1+t^2}$: $F(t) = C -
\arctan t$, and $t \to \infty$ forces $C = \frac\pi2$: $F(t) =
\frac\pi2 - \arctan t$ on $\intoo0{+\infty}$.

\textbf{7.} Integrate by parts on $[A, R]$ with $\sin x =
(-\cos x)'$ and let $R \to \infty$:
\[
\int_A^{\infty}\eu^{-tx}\frac{\sin x}x\dd x
= \frac{\cos A\;\eu^{-tA}}{A}
- \int_A^\infty \cos x\,\Bigl(\frac tx +
\frac1{x^2}\Bigr)\eu^{-tx}\dd x .
\]
Bounding $\abs{\cos} \leq 1$: the first term is $\leq \frac1A$;
the integral is at most $\int_A^\infty t\eu^{-tx}\frac{\dd x}x
+ \int_A^\infty\frac{\dd x}{x^2} \leq \frac1A\int_0^\infty
t\eu^{-tx}\dd x + \frac1A = \frac2A$. Total: $\leq \frac 3A$,
uniformly for $t \in [0, 1]$ (the case $t = 0$ included). Now
\[
\abs{F(t) - D} \leq
\Bigl|\int_0^A(\eu^{-tx} - 1)\,\frac{\sin x}x\,\dd x\Bigr| +
\frac6A .
\]
On $[0, A]$: $\abs{\eu^{-tx} - 1} \leq tx \leq tA$ and
$\abs{\frac{\sin x}x} \leq 1$, so the first term is at most
$tA^2$. Choose $A$ with $\frac6A < \varepsilon$, then $t <
\varepsilon/A^2$: $\abs{F(t) - D} < 2\varepsilon$.

\textbf{8.} Therefore $D = \lim_{t\to0^+}F(t) =
\lim_{t\to0^+}\bigl(\frac\pi2 - \arctan t\bigr) =
\frac\pi2$.

\textbf{9.} By parts ($u = \sin^2x$, $v' = x^{-2}$):
\[
\int_0^\infty\frac{\sin^2x}{x^2}\dd x
= \Bigl[-\frac{\sin^2x}{x}\Bigr]_0^\infty +
\int_0^\infty\frac{2\sin x\cos x}{x}\dd x
= \int_0^\infty\frac{\sin 2x}{x}\dd x = D = \frac\pi2
\]
(substitute $u = 2x$ in the last step; boundary terms vanish:
$\sin^2 x/x \to 0$ at both ends).

\textbf{10.} By parts ($u = 1 - \cos x$, $v' = x^{-2}$):
$\int_0^\infty\frac{1 - \cos x}{x^2}\dd x =
\int_0^\infty\frac{\sin x}x\dd x = \frac\pi2$. Consistency:
$1 - \cos x = 2\sin^2\frac x2$, and the substitution $x = 2u$
turns $\int\frac{2\sin^2(x/2)}{x^2}\dd x$ into
$\int\frac{\sin^2u}{u^2}\dd u$: the two computations agree.

\textbf{11.} $\int_0^\infty\frac{\sin(ax)}x\dd x = \frac\pi2$
for every $a > 0$ (substitute $u = ax$: the integral is
scale-invariant); $\int_\R\eu^{-ax^2}\dd x = \sqrt{\pi/a}$
(substitute $u = \sqrt a\,x$). Scaling in the argument leaves
the Dirichlet integral fixed and divides the Gaussian by
$\sqrt a$.

\textbf{12.} A dominator valid for all $t \in [0,1]$ must
dominate $\sup_{t\in[0,1]}\abs{\eu^{-tx}\frac{\sin
x}x} = \abs{\frac{\sin x}x}$, which is not integrable
(\cref{exo:b3:lebesgue:6}): DCT cannot cross $t = 0$. The
substitute of question 7 --- tails uniformly small in the
parameter, compact part handled by DCT --- is the standard
``uniform integrability'' pattern, and reappears whenever
conditional convergence meets limit interchange.
\end{solution}
